% page 264 du fac-simile (image p-263 du scan)
% bloc 165.1 x 246.3 mm ; marge gauche 8.0 mm
\pgc{-12.700mm}{0.000mm}{
\l{\cel{6}256}
\l{}
\l{\sou{kalvario}. (\sou{kristanismo}). I. Kolino quan acensis Iesu Kristo,}
\l{\cel{3}kande lu portis la kruco, e sur la somito di qua lu kruc-}
\l{\cel{3}agesis. - II. Kolino, sur la suprajo di qua esas kruco,}
\l{\cel{3}adube onu iras pilgrime, haltante e pregante ye la loki}
\l{\cel{3}di la voyo, qui memorigas la cirkonstanci precipua di la}
\l{\cel{3}morto di Iesu Kristo. - (\sou{anke metaf.})\cel{2}DEFIS.}
\l{}
\l{\sou{kalzo.}\cel{2}Texuro ek lano, kotono,silko, +nilono, fitigita a la}
\l{\cel{3}pedo ed a lagambo por kovrar olu. - FIS.}
\l{}
\l{\sou{kam}.\cel{2}Partikulo uzata kom termino duesma en la gradi komparala:}
\l{\cel{3}tam... kam ; sama... kam ; tale... kam ; preferar...kam, e c.}
\l{}
\l{\sou{kamo}.\cel{2}(\sou{tekn.})\cel{2}Saliajo kurvo-linea, muntita sur la perimetro}
\l{\cel{3}di rotaco-axo, qua, renkontrante stango fixigita an la pis-}
\l{\cel{3}tono-stango di pistilo, od an la extremajo di pisto-martel-}
\l{\cel{3}ago, sublevas la pistilo, la martelego, e lasas lu retrofalar,}
\l{\cel{3}tale transformante la movo cirklatra kontinua di la transmis-}
\l{\cel{3}arboro aden movo intermitoza, cirklatra o rekta-linea.DEFIS.}
\l{}
\l{\sou{kamalio}.\cel{2}I. En la Mezepoko, protektilo, konsistanta ye fera}
\l{\cel{3}kaloto en qua esas muntita texuro ek mashi qua protektis la}
\l{\cel{3}kolo e la shultri. - II. (\sou{religio katol.})\cel{2}Quaza kapuco pe-}
\l{\cel{3}lerina,quan la sacerdoti, la kantori metas sur la rocheto ;}
\l{\cel{3}mikra mantelo quan metas sur la rocheto la kardinali, la}
\l{\cel{3}episkopi, la prelati, e c. - III. Mikra mantelo mulierala,}
\l{\cel{3}sen maniki, kun o sen kapuco, diversa-forma (segun la modo).}
\l{\cel{3}- DEF.}
\l{}
\l{\sou{kamarado}. La persono qua habitas normale en la sama loko kam ulu,}
\l{\cel{3}e tale divenas, ad ica, kelke familiara. - DEFIRS.}
\l{}
\l{\sou{kamarilio}.\cel{2}Ensemblo ek la personi qui vivas intime kun princo,}
\l{\cel{3}e probas influar lu politikale. (Vorto desprizala).-DEFIS.}
\l{}
\l{\sou{kambiar}.\cel{2}(\sou{trans.,}\cel{1}po). Donar ad ulu (ulo) e recevar de lu}
\l{\cel{3}(altra kozo kom equivaloza). - IS.}
\l{}
\l{\sou{kambrar}.\cel{2}(\sou{trans.)}\cel{3}Kurvigar arkatre. - EF.}
\l{}
\l{\sou{kame}o. Lapido (onixo, sardonixo, e c.) skultita ad-reliefe,}
\l{\cel{3}ube onu uzas la kolori diversa di la strati por nuancizar}
\l{\cel{3}la parti diversa di la laboruro. - DEFIRS.}
\l{}
\l{\sou{kamelo}. (\sou{zool.)}\cel{2}Quadripedo ruminera, longa-gamba, kun un o}
\l{\cel{3}plura gibi sur la dorso, quan la Arabi, la Orientani, uzas}
\l{\cel{3}kom charjo-bestio. - L.\cel{1}\sou{camelus bactrianus}. - DEFIS.}
\l{}
\l{\sou{kameleono}. I. (\sou{zool.}) Reptero sauria, quaza lacerto kun grosa}
\l{\cel{3}kapo, pelo shagrinatra, qua divenas diversa-kolora - ne}
\l{\cel{3}(quankam onu kredis lo) pro reflektar la objekti vicina,}
\l{\cel{3}ma segun ke la korpo esas plu o min plena ye aero. - L.}
\l{\cel{3}\sou{chameleo vulga}ris.- II. (\sou{metaf.)}\cel{2}Persono qua modifikas}
\l{\cel{3}sua opiniono, sua konduto, segun la okazioni, segun la cir-}
\l{\cel{3}konstanci. - DEFIRS.}
}
